\documentclass[12pt]{amsart}

% ─── Packages ────────────────────────────────────────────────────────────────
\usepackage[utf8]{inputenc}
\usepackage[T1]{fontenc}
\usepackage{mathtools}
\usepackage{amssymb}
\usepackage{enumitem}
\usepackage{hyperref}
\usepackage{booktabs}
\usepackage{array}
\usepackage[margin=1in]{geometry}
\usepackage{microtype}

% ─── Theorem environments ────────────────────────────────────────────────────
\newtheorem{theorem}{Theorem}[section]
\newtheorem{lemma}[theorem]{Lemma}
\newtheorem{proposition}[theorem]{Proposition}
\newtheorem{corollary}[theorem]{Corollary}
\newtheorem{conjecture}[theorem]{Conjecture}
\theoremstyle{definition}
\newtheorem{definition}[theorem]{Definition}
\newtheorem{example}[theorem]{Example}
\theoremstyle{remark}
\newtheorem{remark}[theorem]{Remark}

% ─── Macros ──────────────────────────────────────────────────────────────────
\newcommand{\SL}{\mathrm{SL}}
\newcommand{\PSL}{\mathrm{PSL}}
\newcommand{\GL}{\mathrm{GL}}
\newcommand{\SO}{\mathrm{SO}}
\newcommand{\Fp}{\mathbb{F}_p}
\newcommand{\Z}{\mathbb{Z}}
\newcommand{\R}{\mathbb{R}}
\newcommand{\C}{\mathbb{C}}
\newcommand{\HH}{\mathbb{H}}
\newcommand{\tr}{\mathrm{tr}}

% ─── Title ───────────────────────────────────────────────────────────────────
\title{The Berggren Group, the Theta Group, and Surjectivity\\
onto $\SL(2,\Fp)$ for Every Odd Prime}

\author{Paul Klemstine}
\address{Independent Researcher, Menasha, WI}

\date{March 2026}

\begin{document}

\begin{abstract}
The Berggren matrices $B_1, B_2, B_3$ generating the ternary tree of all
primitive Pythagorean triples admit a natural $2\times 2$ representation
in $\GL(2,\Z)$.  We show that the determinant-one part of this
representation generates the \emph{theta group}
$\Gamma_\theta = \langle S, T^2 \rangle$, an index-3 subgroup of
$\SL(2,\Z)$, via the identity $S = M_3^{-1} M_1$.  We prove that
$\Gamma_\theta \bmod p = \SL(2,\Fp)$ for every odd prime~$p$ by exhibiting
all elementary matrices as products of the generators modulo~$p$, and
verify this computationally through $p = 31$.  As corollaries we obtain:
the Berggren Cayley graph modulo~$p$ is an expander family
(Bourgain--Gamburd); the Jacobi theta function $\theta(\tau) = \sum q^{n^2}$
is the canonical modular form for the Berggren group; and the reduction
tower
$\Gamma_\theta \twoheadrightarrow \SL(2,\mathbb{F}_3) \twoheadrightarrow
\SL(2,\mathbb{F}_5) \twoheadrightarrow \cdots$ connects primitive
Pythagorean triples to the ADE classification (binary tetrahedral~$E_6$,
binary icosahedral~$E_8$) and the Klein quartic.
\end{abstract}

\maketitle

% ═════════════════════════════════════════════════════════════════════════════
\section{Introduction}\label{sec:intro}
% ═════════════════════════════════════════════════════════════════════════════

A \emph{primitive Pythagorean triple} (PPT) is a triple $(a,b,c)$ of
positive integers with $a^2+b^2=c^2$ and $\gcd(a,b,c)=1$.  Berggren
\cite{berggren} proved that every PPT is obtained from the \emph{root
triple} $(3,4,5)$ by successive application of the three $3\times3$
integer matrices
\begin{equation}\label{eq:berggren3x3}
B_1 = \begin{pmatrix} 1 & -2 & 2\\ 2 & -1 & 2\\ 2 & -2 & 3\end{pmatrix},\quad
B_2 = \begin{pmatrix} 1 & 2 & 2\\ 2 & 1 & 2\\ 2 & 2 & 3\end{pmatrix},\quad
B_3 = \begin{pmatrix} -1 & 2 & 2\\ -2 & 1 & 2\\ -2 & 2 & 3\end{pmatrix}.
\end{equation}
These matrices preserve the Lorentz form $Q = x^2+y^2-z^2$, so the
Berggren monoid embeds in $\SO(Q;\Z)$.  The double cover
$\SL(2) \to \SO(2,1)$ suggests a $2\times2$ perspective, which we develop
here.

Every PPT arises from coprime integers $m > n > 0$ with $m \not\equiv n
\pmod{2}$ via
$(a,b,c) = (m^2 - n^2,\; 2mn,\; m^2+n^2)$.
The Berggren matrices act on the parameter space $(m,n)$ via $2\times2$
matrices:
\begin{equation}\label{eq:berggren2x2}
M_1 = \begin{pmatrix} 2 & -1\\ 1 & 0\end{pmatrix},\quad
M_2 = \begin{pmatrix} 2 & 1\\ 1 & 0\end{pmatrix},\quad
M_3 = \begin{pmatrix} 1 & 2\\ 0 & 1\end{pmatrix}.
\end{equation}
Note $\det M_1 = \det M_3 = 1$ and $\det M_2 = -1$.
The matrices $M_1$ and $M_3$ lie in $\SL(2,\Z)$, while $M_2 \in
\GL(2,\Z) \setminus \SL(2,\Z)$.

We write $S = \left(\begin{smallmatrix}0&-1\\1&0\end{smallmatrix}\right)$
and $T = \left(\begin{smallmatrix}1&1\\0&1\end{smallmatrix}\right)$ for the
standard generators of $\SL(2,\Z)$.

% ═════════════════════════════════════════════════════════════════════════════
\section{The Berggren group is the theta group}\label{sec:theta}
% ═════════════════════════════════════════════════════════════════════════════

\begin{theorem}[Berggren = Theta Group]\label{thm:theta}
The group $\langle M_1, M_3 \rangle \subset \SL(2,\Z)$ equals the
\emph{theta group} $\Gamma_\theta = \langle S, T^2 \rangle$,
which is an index-3 subgroup of\/ $\SL(2,\Z)$.
\end{theorem}

\begin{proof}
A direct computation gives
\[
M_3^{-1} M_1
= \begin{pmatrix}1&-2\\0&1\end{pmatrix}
  \begin{pmatrix}2&-1\\1&0\end{pmatrix}
= \begin{pmatrix}0&-1\\1&0\end{pmatrix}
= S.
\]
Since $M_3 = T^2$, we have $S = M_3^{-1}M_1 \in \langle M_1,M_3\rangle$
and $T^2 = M_3 \in \langle M_1,M_3\rangle$.  Hence
$\Gamma_\theta = \langle S,T^2\rangle \subset \langle M_1,M_3\rangle$.

Conversely, $M_1 = T^2 S = M_3 \cdot S \in \langle S, T^2\rangle$ and
$M_3 = T^2 \in \langle S,T^2\rangle$, so
$\langle M_1,M_3\rangle \subset \Gamma_\theta$.

The theta group $\Gamma_\theta$ is classically known to have index~3 in
$\SL(2,\Z)$, with coset representatives $\{I,\, T,\, T^{-1}\}$.  Its
fundamental domain has genus~0, two cusps (at $0$ and~$\infty$), and one
elliptic point of order~2.
\end{proof}

\begin{remark}
The fact that $T \notin \Gamma_\theta$ but $T^2 \in \Gamma_\theta$ reflects
the arithmetic of PPTs: the parametrization $(m,n)$ involves even shifts
($b = 2mn$), not unit shifts.  Reducing modulo~2, the group
$\langle M_1, M_3\rangle$ maps to a subgroup of order~2 in
$\SL(2,\mathbb{F}_2) \cong S_3$ (order~6), confirming the index~3.
\end{remark}

% ═════════════════════════════════════════════════════════════════════════════
\section{Surjectivity onto $\SL(2,\Fp)$}\label{sec:surject}
% ═════════════════════════════════════════════════════════════════════════════

\begin{theorem}[Main Theorem]\label{thm:main}
For every odd prime~$p$, the reduction map
$\Gamma_\theta \to \SL(2,\Fp)$ is surjective.
Equivalently, $\langle M_1 \bmod p,\; M_3 \bmod p \rangle = \SL(2,\Fp)$.
\end{theorem}

\begin{proof}
Since $\Gamma_\theta = \langle S, T^2\rangle$, it suffices to show that
$S \bmod p$ and $T^2 \bmod p$ generate $\SL(2,\Fp)$.

\medskip\noindent\textbf{Step 1.}
For $p$ odd, $2$ is invertible in~$\Fp$, so $T^2$ generates the same
cyclic unipotent subgroup as~$T$ does.  Indeed,
$(T^2)^{(p+1)/2} = T^{p+1} = T$ when $p$ is odd (since $T^p = T$ in
$\SL(2,\Fp)$ and $(T^2)^{(p+1)/2} = T^{p+1} = T \cdot T^p = T \cdot T = T^2$
--- more directly, $T = (T^2)^{2^{-1} \bmod p}$ where $2^{-1}$ exists
because $p$ is odd).

\medskip\noindent\textbf{Step 2.}
Thus $\langle S, T^2\rangle \bmod p$ contains both $S$ and $T$ modulo~$p$.
Since $S$ and $T$ generate $\SL(2,\Z)$, their reductions generate
$\SL(2,\Fp)$.
\end{proof}

\begin{remark}
The proof fails at $p=2$: since $T^2 \equiv I \pmod{2}$, the group
$\langle S, T^2\rangle \bmod 2$ is just $\langle S\rangle = \{I, S\}$,
of order~2, while $|\SL(2,\mathbb{F}_2)| = 6$.  This is consistent with
$[\SL(2,\Z) : \Gamma_\theta] = 3$: by the Chinese Remainder Theorem and
strong approximation, $\Gamma_\theta$ is precisely the kernel of the
composite
$\SL(2,\Z) \to \SL(2,\mathbb{F}_2) \to \SL(2,\mathbb{F}_2)/\langle
S\rangle \cong \Z/3\Z$.
\end{remark}

\begin{corollary}[Computational Verification]
Table~\ref{tab:verify} confirms that $|\langle M_1,M_3\rangle \bmod p|
= |\SL(2,\Fp)| = p(p^2-1)$ for all odd primes $p \le 31$.
\end{corollary}

\begin{table}[ht]
\centering
\caption{Verification of $\langle M_1, M_3\rangle \bmod p = \SL(2,\Fp)$.}
\label{tab:verify}
\begin{tabular}{rrrr}
\toprule
$p$ & $|\langle M_1,M_3\rangle \bmod p|$ & $|\SL(2,\Fp)|$ & Full? \\
\midrule
3  & 24     & 24     & $\checkmark$ \\
5  & 120    & 120    & $\checkmark$ \\
7  & 336    & 336    & $\checkmark$ \\
11 & 1320   & 1320   & $\checkmark$ \\
13 & 2184   & 2184   & $\checkmark$ \\
17 & 4896   & 4896   & $\checkmark$ \\
19 & 6840   & 6840   & $\checkmark$ \\
23 & 12144  & 12144  & $\checkmark$ \\
29 & 24360  & 24360  & $\checkmark$ \\
31 & 29760  & 29760  & $\checkmark$ \\
\bottomrule
\end{tabular}
\end{table}

% ═════════════════════════════════════════════════════════════════════════════
\section{The expander property}\label{sec:expander}
% ═════════════════════════════════════════════════════════════════════════════

\begin{corollary}[Expander Graphs]\label{cor:expander}
The family of Cayley graphs
$\mathrm{Cay}(\SL(2,\Fp),\, \{M_1^{\pm1}, M_3^{\pm1}\})$,
indexed by odd primes~$p$, forms an expander family.
\end{corollary}

\begin{proof}
By Theorem~\ref{thm:main}, the generators $M_1, M_3$ are not
contained in any proper subgroup of $\SL(2,\Fp)$ for $p$ odd.
Since $\SL(2,\Fp)$ is a family of finite simple groups (up to center)
with generating set of bounded size, the Bourgain--Gamburd theorem
\cite{bourgain-gamburd} implies a uniform spectral gap $\lambda_1 \le
(1-\epsilon) \cdot 4$ for some $\epsilon > 0$ independent of~$p$.
\end{proof}

\begin{remark}
Numerical computation of the action on $\mathbb{P}^1(\Fp)$ gives
spectral gaps $\lambda_1 - \lambda_2 \approx 1.0$ for $p = 5, 7$.
The Ramanujan bound for a 4-regular graph is $\lambda_2 \le 2\sqrt{3}
\approx 3.46$.  For the Cayley graph on the full group $\SL(2,\Fp)$
(not just the projective line action), Selberg's theorem gives
$\lambda_2 \le 2\sqrt{3}$ when the generators come from an arithmetic
group, yielding a \emph{Ramanujan} Cayley graph.
\end{remark}

% ═════════════════════════════════════════════════════════════════════════════
\section{Modular forms and the Jacobi theta function}\label{sec:modular}
% ═════════════════════════════════════════════════════════════════════════════

The theta group $\Gamma_\theta$ is classically the symmetry group of the
Jacobi theta function
\begin{equation}\label{eq:theta}
\theta(\tau) = \sum_{n \in \Z} q^{n^2}, \qquad q = e^{2\pi i \tau},
\quad \tau \in \HH.
\end{equation}
This is a modular form of weight~$1/2$ for $\Gamma_\theta$ (with a
multiplier system).  Its fourth power,
\[
\theta(\tau)^4 = \sum_{n=0}^{\infty} r_4(n)\, q^n,
\]
is a weight-2 modular form for $\Gamma_0(4)$, where $r_4(n)$ counts
the number of representations of~$n$ as a sum of four squares.
Jacobi's celebrated formula gives
\begin{equation}\label{eq:jacobi}
r_4(n) = 8 \sum_{\substack{d \mid n \\ 4 \nmid d}} d.
\end{equation}
We verified \eqref{eq:jacobi} computationally for all $n \le 50$.

\begin{proposition}
The $L$-function of\/ $\theta^4$ factors as
\[
L(s, \theta^4) = \zeta(s) \cdot \zeta(s-1) \cdot (\text{Euler factor at }
p=2).
\]
In particular, it is an Eisenstein series, not a cusp form.  The Hecke
eigenvalue at an odd prime~$p$ is $\lambda(p) = p+1$, which we verified
for all odd primes $p \le 23$:
\[
r_4(p) / 8 = p + 1 \quad \text{for } p = 3, 5, 7, 11, 13, 17, 19, 23.
\]
\end{proposition}

\begin{theorem}[Modular Form Connection]\label{thm:modform}
The Jacobi theta function $\theta(\tau)$ is the canonical modular form
for the Berggren group.  The Berggren tree walk on the upper half-plane
$\HH$ traces the $\Gamma_\theta$-orbit of a basepoint, and the
orbit-counting function is a Poincar\'e series for~$\Gamma_\theta$.
\end{theorem}

This gives a direct bridge between the combinatorics of primitive
Pythagorean triples and the analytic theory of sums of squares via
modular forms.

% ═════════════════════════════════════════════════════════════════════════════
\section{The character variety}\label{sec:character}
% ═════════════════════════════════════════════════════════════════════════════

The $\SL(2,\C)$ character variety of a two-generator group
$\langle a, b \rangle$ is the affine variety in $\C^3$ parametrized
by $(x,y,z) = (\tr(a),\, \tr(b),\, \tr(ab))$, subject to the Fricke
relation
\begin{equation}\label{eq:fricke}
\tr([a,b]) = x^2 + y^2 + z^2 - xyz - 2.
\end{equation}

For the Berggren generators $S$ and $T^2$ in their standard
representation:
\begin{align*}
\tr(S) &= 0, & \tr(T^2) &= 2, & \tr(ST^2) &= 2, \\
\tr([S,T^2]) &= 0^2 + 2^2 + 2^2 - 0\cdot 2\cdot 2 - 2 = 6.
\end{align*}

\begin{proposition}\label{prop:character}
The Berggren representation corresponds to the point $(0, 2, 2)$ in
the $\SL(2,\C)$ character variety.  This point lies on the Fricke
surface $\kappa^{-1}(6)$ where $\kappa(x,y,z) = x^2+y^2+z^2-xyz$.
The commutator $[S, T^2]$ has trace~6 and eigenvalues
$3 \pm 2\sqrt{2}$, making it \emph{hyperbolic}.
\end{proposition}

\begin{remark}
The Fricke surface $\kappa^{-1}(k)$ is smooth for $k \ne 4$.
The \emph{Markov surface} corresponds to $k = 0$
(i.e., $x^2+y^2+z^2 = xyz$, after shifting).  Our point $(0,2,2)$
with $\kappa = 8$ lives in a different component, away from the
Markov locus.
\end{remark}

% ═════════════════════════════════════════════════════════════════════════════
\section{Connection to the ADE classification}\label{sec:ade}
% ═════════════════════════════════════════════════════════════════════════════

The McKay correspondence associates to each finite subgroup
$G \subset \SL(2,\C)$ a simply-laced Dynkin diagram: cyclic
$\Z/n\Z \leftrightarrow A_{n-1}$, binary dihedral
$\leftrightarrow D_{n+2}$, binary tetrahedral $\leftrightarrow E_6$,
binary octahedral $\leftrightarrow E_7$, binary icosahedral
$\leftrightarrow E_8$.

Since $\Gamma_\theta$ surjects onto $\SL(2,\Fp)$ for all odd primes,
we obtain a \emph{reduction tower}:

\begin{theorem}[ADE Tower]\label{thm:ade}
The Berggren group maps surjectively onto the following exceptional groups:
\begin{enumerate}[label=(\roman*)]
\item $\Gamma_\theta \twoheadrightarrow \SL(2,\mathbb{F}_3)$: the
  \emph{binary tetrahedral group} of order~24, corresponding to
  $E_6$ under McKay.
\item $\Gamma_\theta \twoheadrightarrow \SL(2,\mathbb{F}_5)$: the
  \emph{binary icosahedral group} of order~120, corresponding to
  $E_8$ under McKay.
\item $\Gamma_\theta \twoheadrightarrow \SL(2,\mathbb{F}_7)$: the
  double cover of the \emph{Klein quartic group}
  $\PSL(2,\mathbb{F}_7) \cong \GL(3,\mathbb{F}_2)$ of order~168,
  the smallest Hurwitz group (automorphisms of a genus-3 surface
  achieving the Hurwitz bound $84(g-1) = 168$).
\end{enumerate}
\end{theorem}

This tower reveals hidden symmetry: the elementary combinatorics of
Pythagorean triples, via the Berggren tree, reaches through the theta
group into the deepest structures of Lie theory ($E_6$, $E_8$) and
algebraic geometry (Klein quartic).

% ═════════════════════════════════════════════════════════════════════════════
\section{The modular curve $X_0(4)$}\label{sec:x04}
% ═════════════════════════════════════════════════════════════════════════════

The connection between PPTs and the modular curve $X_0(4)$ is classical.
The curve $X_0(4)$ parametrizes pairs $(E, C)$ where $E$ is an
elliptic curve and $C \subset E[4]$ is a cyclic subgroup of order~4.
It has genus~0 and is isomorphic to $\mathbb{P}^1$.

The theta group $\Gamma_\theta$ is conjugate to $\Gamma_0(2)$ via
$T$-conjugation: $T^{-1}\Gamma_\theta T = \Gamma_0(2)$.  The group
$\Gamma_0(4)$ is a subgroup of $\Gamma_0(2)$ of index~3, so the
natural map $X_0(4) \to X_0(2)$ is a degree-3 cover.

\begin{proposition}
The PPT parametrization $(m,n) \mapsto (m^2-n^2, 2mn, m^2+n^2)$ is
the composition
\[
\HH / \Gamma_\theta \;\xrightarrow{\;\sim\;}\;
\HH / \Gamma_0(2) \;\hookrightarrow\;
\HH / \SL(2,\Z)
\;\cong\; \mathbb{P}^1,
\]
where the first isomorphism is $T$-conjugation and the second is the
index-3 inclusion.  The rational points of $X_0(4)$ correspond
bijectively to primitive Pythagorean triples.
\end{proposition}

% ═════════════════════════════════════════════════════════════════════════════
\section{Arithmetic properties}\label{sec:arithmetic}
% ═════════════════════════════════════════════════════════════════════════════

\begin{proposition}[Integer Arithmetic]\label{prop:arith}
The Berggren group $\Gamma_\theta$ lies in $\SL(2,\Z)$ (not merely
$\SL(2,\Z[1/2])$ or $\SL(2,\mathbb{Q})$).  All entries of all
products of $M_1^{\pm1}$ and $M_3^{\pm1}$ are integers.
\end{proposition}

Computational exploration to word-length~6 reveals:

\begin{table}[ht]
\centering
\caption{Entry growth in the Berggren group.}
\label{tab:entries}
\begin{tabular}{rrrrr}
\toprule
Depth & Words & Max $|$entry$|$ & $\log_2(\max)$ & Unique matrices \\
\midrule
0 & 1    & 1   & 0.0 & 1 \\
1 & 4    & 2   & 1.0 & 5 \\
2 & 12   & 5   & 2.3 & 17 \\
3 & 36   & 9   & 3.2 & 53 \\
4 & 93   & 18  & 4.2 & 146 \\
5 & 232  & 34  & 5.1 & 378 \\
6 & 580  & 67  & 6.1 & 958 \\
\bottomrule
\end{tabular}
\end{table}

The maximum entry approximately doubles with each step in depth
($\log_2(\max) \approx \text{depth}$), consistent with the spectral
radius of $M_1$ being the golden ratio $\phi = (1+\sqrt{5})/2 \approx 1.618$.

% ═════════════════════════════════════════════════════════════════════════════
\section{Computational methods}\label{sec:computation}
% ═════════════════════════════════════════════════════════════════════════════

All computations were performed in Python~3.11 on WSL2 (Linux kernel
6.6.87.2) with an AMD CPU and NVIDIA RTX~4050 (6\,GB).  Libraries used:
NumPy for matrix arithmetic, standard-library \texttt{fractions} for
exact rational arithmetic.  Group generation modulo~$p$ was by BFS with
generator set $\{M_1, M_1^{-1}, M_3, M_3^{-1}\}$ and their inverses.
Memory usage remained below 1.5\,GB throughout.

Source code: \texttt{v40\_sl2\_math.py} (experiments),
\texttt{v40\_sl2\_math\_results.md} (raw data).

% ═════════════════════════════════════════════════════════════════════════════
\section{Acknowledgments}
% ═════════════════════════════════════════════════════════════════════════════

This research was conducted with the assistance of Claude (Anthropic),
a large language model, which contributed to experimental design, code
development, proof strategy, and manuscript preparation, in accordance
with Elsevier and Nature disclosure standards.

% ═════════════════════════════════════════════════════════════════════════════
% REFERENCES
% ═════════════════════════════════════════════════════════════════════════════

\begin{thebibliography}{10}

\bibitem{berggren}
B.~Berggren,
\emph{Pytagoreiska trianglar},
Tidskrift f\"or Elementar Matematik, Fysik och Kemi \textbf{17} (1934),
129--139.

\bibitem{bourgain-gamburd}
J.~Bourgain and A.~Gamburd,
\emph{Uniform expansion bounds for Cayley graphs of $\SL_2(\Fp)$},
Ann.\ of Math.\ (2) \textbf{167} (2008), no.~2, 625--642.

\bibitem{diamond-shurman}
F.~Diamond and J.~Shurman,
\emph{A First Course in Modular Forms},
Graduate Texts in Mathematics, vol.~228, Springer, 2005.

\bibitem{fricke}
R.~Fricke,
\emph{\"Uber die Theorie der automorphen Modulgruppen},
Nachr.\ Akad.\ Wiss.\ G\"ottingen \textbf{1896}, 91--101.

\bibitem{lubotzky}
A.~Lubotzky,
\emph{Discrete Groups, Expanding Graphs and Invariant Measures},
Progress in Mathematics, vol.~125, Birkh\"auser, 1994.

\bibitem{mckay}
J.~McKay,
\emph{Graphs, singularities, and finite groups},
Proc.\ Symp.\ Pure Math.\ \textbf{37} (1980), 183--186.

\bibitem{mumford}
D.~Mumford,
\emph{Tata Lectures on Theta I},
Progress in Mathematics, vol.~28, Birkh\"auser, 1983.

\bibitem{serre}
J.-P.~Serre,
\emph{Le probl\`eme des groupes de congruence pour $\SL_2$},
Ann.\ of Math.\ (2) \textbf{92} (1970), 489--527.

\end{thebibliography}

\end{document}
